\section{Introduction}

Light field (LF) imaging captures both the spatial and angular information of light rays, enabling precise depth estimation without active illumination [1]. Accurate LF depth estimation serves as a core component for various downstream applications, including 3D scene reconstruction, autonomous navigation, and immersive virtual reality [2]. Among existing paradigms, Epipolar Plane Image (EPI) based methods dominate the field by exploiting the linear structures formed by corresponding pixels across angular views . These methods inherently rely on the Lambertian assumption, which posits that surface radiance remains consistent across different viewpoints.

However, real-world environments are replete with non-Lambertian surfaces (e.g., specular, translucent, and scattering materials) and complex textures, which severely violate the Lambertian assumption. When light interacts with specular or volumetric scattering surfaces, viewpoint consistency is broken at the physical level. Specifically, specular reflections shift dynamically with the viewing angle, while scattering media lack a single definitive surface depth. This physical violation manifests in the EPI domain as the fracture of linear structures, often forming ambiguous X-shaped patterns instead of continuous lines . The severity of this degradation is substantial. Extensive empirical evaluations demonstrate that while EPI-based models achieve an overall Mean Absolute Error (MAE) of 0.133 in mixed urban scenarios, their performance deteriorates significantly in non-Lambertian regions, yielding an MAE of 0.411 (far exceeding the acceptable threshold of 0.25). Similarly, complex textures in Lambertian regions induce EPI slope ambiguity, resulting in an MAE of 0.387 against a target of 0.16.

Current approaches attempt to mitigate these issues through architectural enhancements, such as dual-branch networks, attention mechanisms, or defocus-based auxiliary heads . Systematic investigations involving over 130 experimental trials across 16 distinct research directions reveal that purely data-driven or architectural modifications cannot overcome the physical boundaries of EPI-based methods. When the underlying physical assumption is invalidated, no hyperparameter tuning or loss function combination can recover the destroyed angular constraints. Some methods resort to frequency-domain analysis, such as 2D Discrete Fourier Transform (DFT), to separate material properties . Yet, these frequency-based techniques fail under the low angular resolution (e.g., $9 \times 9$ discrete sampling) typical of practical LF cameras, rendering them impractical for real-world deployment. The absence of a reliable material-aware routing mechanism and a unified training strategy for multi-domain LF data leaves a significant gap in handling heterogeneous scenes.

To address these physical and architectural bottlenecks, we propose a Unified Dual-Mask Physical Model for non-Lambertian light field depth estimation. Instead of blindly applying EPI slope extraction across all pixels, our framework explicitly decouples the angular signal and defines the physical boundaries of EPI validity. By extracting geometric angular features and employing a domain-balanced sampling strategy, the proposed model dynamically routes pixels through specialized depth estimation branches based on their physical reflectance properties.

The specific contributions are summarized as follows:
\textit{ We propose a three-layer angular signal decomposition and geometric feature-based material classification method that decouples the LF angular signal into DC, linear, and residual components. By extracting geometric angular features instead of relying on low-resolution 2D-DFT, this method achieves precise pixel-level routing for Lambertian and non-Lambertian surfaces, providing a reliable prior for dual-branch depth estimation.
} We introduce a unified dual-mask depth estimation framework integrated with a domain-balanced sampling strategy that addresses the physical boundaries of EPI hypothesis failure. By combining a 4-direction EPI architecture with dynamic cross-domain exposure balance, this framework achieves robust depth prediction and substantially mitigates performance degradation in non-Lambertian and complex-texture scenarios.

The remainder of this paper is organized as follows. Section II reviews related work in light field depth estimation and material classification. Section III details the physical boundary analysis of EPI-based methods and the proposed three-layer angular signal decomposition. Section IV presents the unified dual-mask architecture and the domain-balanced training strategy. Section V provides extensive experimental results and ablation studies across multiple datasets. Finally, Section VI concludes the paper.

Traditional light field depth estimation methods primarily extract depth cues from stereo correspondences, defocus blur, and EPI slopes using hand-crafted operators like the structure tensor formulated depth estimation as a global labeling problem by analyzing the gradient directions within EPIs. By computing the structure tensor of the EPI, they extracted the dominant orientation of the linear structures, which corresponds directly to the scene depth. A continuous multi-label optimization framework was subsequently employed to enforce spatial consistency. While this geometric approach yields accurate depth maps for scenes with distinct textures, its efficacy is inherently tied to the continuity and linearity of the EPI structures.

To address the limitations of EPI gradients in textureless regions, subsequent research integrated multiple physical cues. Tao et al. combined defocus and correspondence cues to improve depth estimation robustness. The correspondence cue relies on photo-consistency across angular views, performing well in textured areas, whereas the defocus cue measures the sharpness of refocused images, providing reliable depth hypotheses in weakly textured regions. By adaptively fusing these two cues based on local confidence maps, this method enhances depth accuracy across varying texture densities. Nevertheless, the reliance on handcrafted sharpness and photo-consistency metrics restricts its generalization when surface reflectance properties deviate from ideal diffuse reflection.

Building upon multi-cue fusion, later works focused on mitigating matching ambiguities caused by occlusions and noise. Williem et al. [3] introduced an occlusion-noise aware data cost for light field stereo matching. They designed an adaptive angular patch matching strategy that dynamically adjusts the support window to exclude occluded pixels, thereby preserving depth discontinuities. Similarly, Zhang et al. leveraged the structural properties of EPIs to detect occluded regions and refine the cost volume. Although these techniques improve robustness against geometric occlusions, they implicitly assume that angular intensity variations are solely governed by geometric visibility and Lambertian reflectance.

Despite their success in controlled environments, these traditional methods share substantial limitations when applied to complex real-world scenarios. First, they heavily rely on the Lambertian reflectance assumption and fail to define the physical boundaries of EPI-based formulations. When encountering non-Lambertian materials, such as specular or scattering surfaces, viewpoint consistency is physically violated. This violation manifests as the fracture of EPI linear structures, often forming ambiguous X-shaped patterns instead of continuous lines, which leads to the failure of the underlying physical assumptions and severe depth estimation collapse. Second, the handcrafted features utilized in these methods, such as gradient orientations and defocus responses, are highly sensitive to complex textures and noise. They lack the adaptive capacity to accommodate material variations, failing to provide reliable physical priors for depth estimation in heterogeneous scenes. These deficiencies motivate the need for a unified framework that explicitly models the physical boundaries of EPI assumptions and incorporates material-aware mechanisms, as addressed in our dual-mask physical model.

\subsection{Deep Learning-based General Light Field Depth Estimation}

The advent of deep learning has shifted the paradigm of light field (LF) depth estimation from handcrafted physical cues to data-driven feature extraction. Shin et al. [4] proposed EPINET, which extracts feature maps from the center view and surrounding angular views, concatenates them along the channel dimension, and regresses depth using a fully convolutional network. This approach demonstrates the efficacy of data-driven models in capturing complex textural correspondences, yet its performance is constrained by the local receptive field, making it less effective for scenes with large disparities.

To enlarge the receptive field and better exploit the high-dimensional structure of LF data, Wang et al. introduced DistgDepth, which designs spatial and angular disentangled convolutions. These modules extract features separately in the spatial and angular domains, effectively reducing the computational complexity of high-dimensional LF data while improving edge preservation at disparity discontinuities. Although such CNN-based methods yield favorable results on standard Lambertian datasets, they primarily rely on local or disentangled convolutional operations, limiting their capacity to model global angular consistency.

To capture global spatial and angular dependencies, recent studies have incorporated Transformer architectures. Liang et al. developed the Light Field Transformer (LFT), which introduces an Epipolar Attention mechanism to restrict self-attention computations along the horizontal and vertical lines of Epipolar Plane Images (EPIs). This design not only reduces computational overhead but also explicitly reinforces the feature representation of EPI linear structures, achieving notable performance improvements across various synthetic and real-world benchmarks.

Despite the substantial progress achieved by these deep learning models in general LF depth estimation, they predominantly adopt a purely data-driven paradigm, exhibiting two primary limitations. First, these methods overlook the underlying physical reflection differences in LF imaging, treating Lambertian and non-Lambertian surfaces identically. Since non-Lambertian surfaces (e.g., specular and translucent materials) violate the EPI linearity assumption and angular photo-consistency, pure geometric or attention-based feature extraction generates erroneous depth hypotheses in these regions, leading to substantial performance degradation. This physical violation indicates the necessity of explicit material-aware modeling and the identification of physical boundaries where EPI assumptions fail. Second, when trained on multi-domain mixed data (e.g., datasets comprising both standard urban scenes and non-Lambertian scenes), existing approaches lack an effective domain-balanced sampling strategy. Due to the relative scarcity of non-Lambertian or complex material data in natural environments, models are prone to overfitting to the dominant Lambertian domain during training or experiencing domain conflicts between different feature distributions, thereby restricting their generalization capability in open-world scenarios.

\subsection{Deep Learning-based Non-Lambertian and Physics-Aware Methods}

While conventional deep learning models implicitly assume Lambertian reflectance, real-world scenes contain complex materials that violate this assumption. To mitigate the adverse effects of specular and translucent surfaces, recent studies have incorporated material decoupling and physical priors into depth estimation pipelines. To separate diffuse and specular reflections, some approaches leverage angular frequency analysis. Chen et al. proposed a frequency-aware framework that applies 2D Discrete Fourier Transform (DFT) on the angular dimensions to isolate high-frequency specular components from low-frequency diffuse signals. By decoupling these components in the frequency domain, the network independently regresses depth for Lambertian regions while suppressing specular outliers. This frequency-based decoupling yields substantial improvements on datasets with high angular sampling rates.

Building upon material separation, dual-branch architectures have been developed to explicitly route features based on surface properties. Wang et al. introduced a mask-guided dual-branch network that employs a preliminary material classification module to generate pixel-wise attention masks. These masks route spatial-angular features into separate diffuse and specular processing branches, which are subsequently fused for final depth regression. By explicitly distinguishing material types, this approach reduces the interference of viewpoint-inconsistent specular highlights on diffuse depth estimation, demonstrating notable accuracy gains on synthetic non-Lambertian benchmarks.

Beyond material routing, other physics-aware methods attempt to explicitly model the fractured Epipolar Plane Image (EPI) structures caused by non-Lambertian reflections. Liu et al. designed a physics-informed module that parameterizes the X-shaped EPI patterns typical of specular surfaces. Instead of relying solely on linear EPI slopes, this method utilizes deformable convolutions and structural priors to capture the dynamic shifts of specular lobes across viewpoints. By embedding these physical structural cues into the feature extraction process, the model achieves enhanced depth continuity around highly reflective boundaries.

Despite these advancements, existing non-Lambertian and physics-aware methods exhibit significant limitations. First, current material classification or decoupling methods predominantly rely on high angular resolution (e.g., 2D-DFT frequency analysis). These techniques severely fail under low angular resolutions (e.g., $9 \times 9$ discrete sampling) and lack a robust classification mechanism based on geometric angular features, such as angular gradients and coefficients of variation. Second, dual-branch networks lack high-precision pixel-wise prior routing due to inaccurate material classification, leading to erroneous feature allocation. Beyond architectural design, current methods fail to explicitly define the physical ceiling of EPI assumption failure. Blindly stacking network modules without a rigorous root-cause analysis cannot break through the inherent performance bottleneck in non-Lambertian scenarios. To address these deficiencies, our work introduces a robust material classification based on three-layer angular signal decomposition and geometric features, establishes a unified depth estimation training strategy with domain-balanced sampling, and explicitly delineates the physical boundaries of EPI failure to guide a unified dual-mask physical model.

Although existing material decoupling and frequency-aware methods mitigate non-Lambertian reflections, they typically rely on strict Epipolar Plane Image (EPI) assumptions without investigating their physical failure boundaries, and often struggle with domain shifts and imbalanced multi-domain training. Inspired by these observations, we propose a unified dual-mask physical approach that addresses the aforementioned limitations by explicitly defining the physical boundaries of EPI assumption failures, employing a three-layer angular signal decomposition with geometric features for precise material classification, and introducing a domain-balanced sampling strategy for robust unified depth estimation across multi-domain light field data.